\documentclass[12pt]{report}
\usepackage[margin=1in]{geometry}
\begin{document}

\begin{titlepage}
\centering
\vspace*{2cm}
{\Huge\bfseries Advances in Document Processing\\[0.5cm]}
{\Large A Comprehensive Survey\\[2cm]}
{\large Dr.\ Katherine Park\\[0.3cm]}
{\normalsize Department of Computer Science\\
University of Technology\\
kpark@university.edu\\[2cm]}
{\large April 2026\\[3cm]}
{\small Technical Report TR-2026-042}
\end{titlepage}

\begin{abstract}
This report provides a comprehensive survey of document processing techniques,
covering optical character recognition, layout analysis, semantic classification,
and content extraction. We examine the evolution from rule-based systems to modern
deep learning approaches, identifying key challenges in handling diverse document
formats. Our analysis covers 150 papers published between 2015 and 2025, revealing
trends toward end-to-end learning systems and multimodal architectures.
\end{abstract}

\chapter{Introduction}
Document processing is the task of extracting structured information from
unstructured document images or digital files. This field encompasses several
sub-tasks including text recognition, layout analysis, and semantic understanding.

\section{Motivation}
The volume of digital documents continues to grow exponentially. Organizations
need automated tools to extract, classify, and index information from PDF files,
scanned documents, and other formats.

\section{Scope}
This survey focuses on digital PDF processing, excluding scanned document OCR.
We examine techniques for preserving content fidelity during conversion to
structured formats such as Markdown, HTML, and XML.

\chapter{Background}
\section{PDF Format}
The Portable Document Format stores content as positioned drawing commands rather
than logical text flow. This design prioritizes pixel-perfect rendering over
content accessibility.

\section{Markdown Format}
Markdown is a lightweight markup language designed for readability. It supports
headings, lists, tables, images, and inline formatting through simple syntax.
\end{document}